\subsection{Expressive-power map}
\label{sec:p14-expressive-table}

\begin{longtable}{@{}p{0.21\textwidth}p{0.27\textwidth}p{0.42\textwidth}@{}}
\caption{Expressive-power map}\\
\toprule
Comparison & Inclusion or equivalence & Witnessing boundary \\
\midrule
\endfirsthead
\caption[]{Expressive-power map (continued)}\\
\toprule
Comparison & Inclusion or equivalence & Witnessing boundary \\
\midrule
\endhead
\midrule\multicolumn{3}{r}{\emph{continued on next page}}\\
\endfoot
\endlastfoot
modal logic vs first-order logic & standard translation embeds modal logic in FO & bisimulation-invariant boundary \cite{VanBenthem1983} \\
CPL vs first-order logic & propositional formulas embed as nullary-predicate cases & domain-sensitive quantification \cite{Enderton2001} \\
FO vs existential SO on finite structures & proper inclusion & ESO captures NP \cite{Fagin1974} \\
FO vs least fixed-point logic on ordered finite structures & proper inclusion & LFP captures P \cite{Immerman1986} \\
LFP vs ESO on ordered finite structures & unconditional inclusion & \(\mathsf P\subseteq\mathsf{NP}\); equality exactly when \(\mathsf P=\mathsf{NP}\) \cite{Immerman1999} \\
LTL vs FO on words & LTL with past captures FO point properties & Kamp theorem \cite{Kamp1968} \\
CTL vs LTL & incomparable & branching/path witnesses \cite{EmersonHalpern1986} \\
CTL and LTL vs CTL* & both proper fragments & unrestricted state/path nesting \cite{EmersonHalpern1986} \\
modal \(\mu\)-calculus vs modal logic & proper extension & fixed-point reachability \cite{Kozen1983} \\
full SO vs Henkin SO & same syntax, different validities & full versus general relation ranges \cite{Henkin1950,Shapiro1991} \\
\bottomrule
\end{longtable}

All comparisons hold over the displayed structure class and semantics. A
translation that preserves satisfiability need not preserve formula size,
validity, proof length, or intended interpretation.
